
\documentclass[fleqn,10pt]{SelfArx} 
\setlength{\columnsep}{0.55cm} 
\setlength{\fboxrule}{0.75pt} 
\definecolor{color1}{RGB}{0,0,90} 
\definecolor{color2}{RGB}{0,20,20} 
\usepackage{hyperref} 

\hypersetup{hidelinks,colorlinks,breaklinks=true,urlcolor=color2,citecolor=color1,linkcolor=color1,bookmarksopen=false,pdftitle={Title},pdfauthor={Author}}

%-------------------------------- Paquetes útiles --------------------------------

\usepackage{graphicx}
\usepackage[utf8]{inputenc}
\usepackage{url}
\usepackage{amsmath}

%\usepackage{subfigure}
%\usepackage{cite}
%\usepackage{stfloats}
%\usepackage{array}

%------------------------------ Título y encabezado ------------------------------


\JournalInfo{Fecha dd/mm/aa (Día de la práctica)} 
\Archive{Subgrupo N. ??? \\
{\color[rgb]{1,0,0} Nota: }[ \:\:\: ]} 

\PaperTitle{TÍTULO DE LA PRÁCTICA \color[rgb]{1,0,0}{(TODO EN MAYÚSCULA)}} 
\Authors{Estudiante 1 (código), Estudiante 2 (código)}

\Keywords{} 
\newcommand{\keywordname}{} 


%----------------------------------------------------------------------------------

\Abstract{Breve descripción de la práctica de laboratorio y la importancia de éste en la física, además de una posible aplicación de los principios y el conocimiento acerca del tema. \color[rgb]{1,0,0}{POR FAVOR BUENA ORTOGRAFÍA, ESTO RESTA CALIDAD AL TRABAJO}}
\begin{document}
\flushbottom % Asigna la misma altura a todas las páginas
\maketitle
\thispagestyle{empty} % Evita la numeración de las páginas

%----------------------------------------------------------------------------------
\section{MARCO TEÓRICO}

\begin{equation}
\cos^3 \theta =\frac{1}{4}\cos\theta+\frac{3}{4}\cos 3\theta
\label{Eq. 1}
\end{equation}

\subsection{Subsección 1}

%----------------------------------------------------------------------------------
\section{OBJETIVOS}
\subsection{Objetivo Principal}

Objetivo principal

\subsection{Objetivos Específicos}
\begin{itemize}[noitemsep]

\item Objetivo específico 1
\item Objetivo específico 2
\item Objetivo específico 3

\end{itemize}

%----------------------------------------------------------------------------------
\section{MATERIALES Y MÉTODOS}

\subsection{Materiales e Instrumentos}
\begin{enumerate}[noitemsep] 

\item Material e instrumentos 1
\item Material e instrumentos 2
\item Material e instrumentos 3

\end{enumerate}


\subsection{Metodología}
\begin{enumerate}[noitemsep] 

\item Paso 1 con explicación
\item Paso 2 con explicación
\item Paso 3 con explicación

\end{enumerate}

%---------------------------------------------------------------------------------------
\section{RESULTADOS Y DISCUSIÓN}

Los datos obtenidos deben estar representados en tablas, estos valores pueden incluir de una vez las desviaciones de cada uno de los datos
Recordar que todas las medidas y resultados deben tener las desviaciones y errores pertinentes, se deben calcular los porcentajes de error  en el caso de comparar con respecto a un valor teórico.
La discusión de los datos debe tener argumentos que sean conexos con los principios físicos plasmados en el marco teóricos.

%---------------------------------------------------------------------------------------
\section{CONCLUSIONES}






%---------------------------------------------------------------------------------------
\bibliographystyle{unsrt}
\bibliography{sample}






\end{document}
